% Adapted from class/easy/wwpd.tex for the dataclass interface.
\question What would Python display? Run these statements in order.
If a statement causes an error, identify the error and continue.

\begin{blocksection}
\begin{lstlisting}
>>> a = Link(1, Link(2, Link(3)))
>>> a.first
\end{lstlisting}
\begin{solution}[.35in]
\begin{lstlisting}
1
\end{lstlisting}
\end{solution}
\end{blocksection}

\begin{blocksection}
\begin{lstlisting}
>>> a.rest.first
\end{lstlisting}
\begin{solution}[.35in]
\begin{lstlisting}
2
\end{lstlisting}
\end{solution}
\end{blocksection}

\begin{blocksection}
\begin{lstlisting}
>>> a.rest.rest.first
\end{lstlisting}
\begin{solution}[.35in]
\begin{lstlisting}
3
\end{lstlisting}
\end{solution}
\end{blocksection}

\begin{blocksection}
\begin{lstlisting}
>>> a.rest.rest.rest
\end{lstlisting}
\begin{solution}[.35in]
\begin{lstlisting}
()
\end{lstlisting}
\end{solution}
\end{blocksection}

\begin{blocksection}
\begin{lstlisting}
>>> isinstance(a.rest, Link)
\end{lstlisting}
\begin{solution}[.35in]
\begin{lstlisting}
True
\end{lstlisting}
\end{solution}
\end{blocksection}

\begin{blocksection}
\begin{lstlisting}
>>> isinstance(a.rest.rest.rest, Link)
\end{lstlisting}
\begin{solution}[.35in]
\begin{lstlisting}
False
\end{lstlisting}
\end{solution}
\end{blocksection}

\begin{blocksection}
\begin{lstlisting}
>>> a.rest.rest.rest.first
\end{lstlisting}
\begin{solution}[.35in]
\begin{lstlisting}
AttributeError: 'tuple' object has no attribute 'first'
\end{lstlisting}
\end{solution}
\end{blocksection}

\begin{blocksection}
\begin{lstlisting}
>>> a.rest.rest
\end{lstlisting}
\begin{solution}[.35in]
\begin{lstlisting}
Link(first=3, rest=())
\end{lstlisting}
\end{solution}
\end{blocksection}

\begin{blocksection}
\begin{lstlisting}
>>> Link(1, Link(2))
\end{lstlisting}
\begin{solution}[.35in]
\begin{lstlisting}
Link(first=1, rest=Link(first=2, rest=()))
\end{lstlisting}
\end{solution}
\end{blocksection}

\begin{blocksection}
\begin{lstlisting}
>>> print(a)
\end{lstlisting}
\begin{solution}[.35in]
\begin{lstlisting}
(1 2 3)
\end{lstlisting}
\end{solution}
\end{blocksection}

\begin{questionmeta}
\textbf{Teaching Tips}\par\nopagebreak[4]
Suggested Time: 6--8 min; Difficulty: Easy
\nopagebreak[4]
\begin{itemize}
    \item \textbf{Follow the structure:} Each \lstinline{rest} skips one element. Name the resulting value before taking another attribute.
    \item \textbf{Recognize the base case:} \lstinline{()} is an empty linked list, not a \lstinline{Link} instance. It has no \lstinline{first} or \lstinline{rest}.
    \item \textbf{Read the representations:} The dataclass representation shows \lstinline{first=} and \lstinline{rest=}; \lstinline{print} uses the supplied formatter to show the sequence as \lstinline{(1 2 3)}. Both describe the same element sequence.
\end{itemize}
\end{questionmeta}
